\documentclass[11pt,a4paper]{article}
\usepackage[margin=25mm]{geometry}
\usepackage{amsmath,amssymb}
\usepackage{fontspec}
\setmainfont{SimSun}
\setsansfont{Microsoft YaHei}
\setmonofont{Consolas}
\XeTeXlinebreaklocale "zh"
\XeTeXlinebreakskip = 0pt plus 1pt
\usepackage{graphicx}
\usepackage{float}
\floatplacement{figure}{H}
\renewcommand{\figurename}{图}
\renewcommand{\tablename}{表}
\setcounter{secnumdepth}{-1}
\usepackage{longtable,booktabs,array}
\usepackage{hyperref}
\hypersetup{colorlinks=true,linkcolor=black,urlcolor=blue}
\setlength{\parindent}{2em}
\setlength{\parskip}{4pt}
\linespread{1.2}
\widowpenalty=10000
\clubpenalty=10000
\setkeys{Gin}{width=\linewidth,keepaspectratio}
\providecommand{\tightlist}{\setlength{\itemsep}{0pt}\setlength{\parskip}{0pt}}
\providecommand{\passthrough}[1]{#1}
\title{独立机制验证实验报告}
\author{}
\date{2026年9月17日}
\begin{document}
\maketitle
\begin{center}冻结协议结果；不与开发集混合\end{center}
\hypertarget{ux8303ux56f4ux4e0eux6570ux636eux5ba1ux8ba1}{%
\section{1
范围与数据审计}\label{ux8303ux56f4ux4e0eux6570ux636eux5ba1ux8ba1}}

全24个BBOB函数，10/20/40维，实例11--13，各3重复；六种内部配置共3,888运行、907,200算法评价。三种预定历史共1,944轨迹、7,776冻结状态，独立取得1,990,656额外标签。完整日志、共同初始化、评价前信息顺序和源指纹通过全量检查。原计划实例6--10未读取。本阶段不调参数、不删除函数、不并入开发集统计，也不倒置开发门槛的失败结论。

\hypertarget{ux9884ux5b9aux4e3bux8981ux6bd4ux8f83}{%
\section{2 预定主要比较}\label{ux9884ux5b9aux4e3bux8981ux6bd4ux8f83}}

估计正负需按指标解释：P1为fixed03的选中-池内失效差减greedy的对应差；闭环P2为PEM-T减固定.1的终点对数误差，负值更好；单步P2为反馈减均值的归一化真实改善，正值更好。24函数配对，9项统一Holm；分层重采样区间未作多重校正，不能替代校正检验。

\par\noindent\begin{minipage}{\linewidth}

\begin{longtable}[]{@{}llllllll@{}}
\caption{9项预定比较完整结果}\tabularnewline
\toprule
维度 & 比较 & 估计 & 区间下限 & 区间上限 & 正方向函数 & 负方向函数 &
Holm p \\
\midrule
\endfirsthead
\toprule
维度 & 比较 & 估计 & 区间下限 & 区间上限 & 正方向函数 & 负方向函数 &
Holm p \\
\midrule
\endhead
10 & 历史差异之差 & 0.0851 & 0.0311 & 0.1234 & 20 & 4 & 0.0015 \\
20 & 历史差异之差 & 0.1225 & 0.0696 & 0.1948 & 21 & 3 & 0.0006 \\
40 & 历史差异之差 & 0.1449 & 0.0941 & 0.2244 & 17 & 7 & 0.0017 \\
10 & 终点误差差 & 0.1668 & 0.0749 & 0.2391 & 22 & 2 & 5.36e-06 \\
20 & 终点误差差 & 0.0998 & 0.0467 & 0.1905 & 19 & 5 & 0.0017 \\
40 & 终点误差差 & 0.0639 & 0.0130 & 0.1123 & 19 & 5 & 0.0017 \\
10 & 单步改善差 & -0.0170 & -0.0431 & -0.0067 & 1 & 18 & 0.0008 \\
20 & 单步改善差 & -0.0158 & -0.0352 & -0.0079 & 0 & 21 & 0.0005 \\
40 & 单步改善差 & -0.0344 & -0.0630 & -0.0100 & 1 & 20 & 0.0005 \\
\bottomrule
\end{longtable}

\end{minipage}\par

\par\noindent\begin{minipage}{\linewidth}

\begin{longtable}[]{@{}llllll@{}}
\caption{主原型闭环终点效果}\tabularnewline
\toprule
维度 & 终点对数差 & 几何误差比 & 误差较大函数 & 误差较小函数 & Holm p \\
\midrule
\endfirsthead
\toprule
维度 & 终点对数差 & 几何误差比 & 误差较大函数 & 误差较小函数 & Holm p \\
\midrule
\endhead
10 & 0.1668 & 1.4681 & 22 & 2 & 5.36e-06 \\
20 & 0.0998 & 1.2584 & 19 & 5 & 0.0017 \\
40 & 0.0639 & 1.1585 & 19 & 5 & 0.0017 \\
\bottomrule
\end{longtable}

\end{minipage}\par

\begin{figure}
\centering
\includegraphics{../results/mechanism_validation_report_v1/endpoint_function_effects.png}
\caption{逐函数终点差。橙色表示PEM-T误差更大，蓝色表示误差更小；每维24函数均保留。}
\end{figure}

\hypertarget{ux5386ux53f2ux4e0eux76d1ux6d4bux5bf9ux8c61}{%
\section{3
历史与监测对象}\label{ux5386ux53f2ux4e0eux76d1ux6d4bux5bf9ux8c61}}

\par\noindent\begin{minipage}{\linewidth}

\begin{longtable}[]{@{}llll@{}}
\caption{T上界的候选池与实际选中失效率}\tabularnewline
\toprule
维度 & 历史 & 候选池失效 & 实际选中失效 \\
\midrule
\endfirsthead
\toprule
维度 & 历史 & 候选池失效 & 实际选中失效 \\
\midrule
\endhead
10 & fixed03 & 0.0479 & 0.1215 \\
10 & greedy & 0.1018 & 0.0903 \\
10 & pool\_random & 0.0859 & 0.0660 \\
20 & fixed03 & 0.0442 & 0.0833 \\
20 & greedy & 0.1494 & 0.0660 \\
20 & pool\_random & 0.0900 & 0.0729 \\
40 & fixed03 & 0.0505 & 0.0799 \\
40 & greedy & 0.1850 & 0.0694 \\
40 & pool\_random & 0.0978 & 0.0694 \\
\bottomrule
\end{longtable}

\end{minipage}\par

\begin{figure}
\centering
\includegraphics{../results/mechanism_validation_report_v1/history_object_differences.png}
\caption{不同历史的实际选中减候选池失效率。方向由结果决定，不预设筛选必然恶化。不同历史并非同池反事实，图中差异不解释为历史策略的独立因果效应。}
\end{figure}

\hypertarget{ux6837ux672cux652fux6301ux4e0eux58f0ux660e}{%
\section{4
样本支持与声明}\label{ux6837ux672cux652fux6301ux4e0eux58f0ux660e}}

\begin{figure}
\centering
\includegraphics{../results/mechanism_validation_report_v1/effective_support.png}
\caption{时间及局部时空上界的状态平均ESS。每轨迹四个实际检查点；首个检查点分别41/100、61/200、101/400。虚线20是工程支持阈值，不是覆盖保证。}
\end{figure}

检查点支持与回退完整表另存CSV。闭环声明汇总表联合保存声明总数、错误数、探索、失效及FDP；无声明FDP为NA。稀疏池内声明表另公开非空状态数、候选数及合并失败比例，不从零错误或低失败比例推出稳定声明成功概率。

\hypertarget{ux89e3ux91caux8fb9ux754cux4e0eux540eux7eedux6574ux5408}{%
\section{5
解释边界与后续整合}\label{ux89e3ux91caux8fb9ux754cux4e0eux540eux7eedux6574ux5408}}

本报告公开全部预定方向、区间与校正检验。需要根据这些实际结果逐项判断原经验命题是否外推，任何反向或接近零结果应进入正文讨论。排序等价性仍为固定规则的代数性质；经验失败不能扩展到所有校准或探索方法。共同状态即期收益不能证明长期因果机制。

与已完成Sobol/CMA-ES参照的共同归一化表仅在本阶段审计后生成；初始误差分母不向外部算法提供。该描述性展示不追加未经协议规定的显著性检验。GP-EI/TuRBO及关键SAEA、应用和独占资源证据仍需补齐。当前报告不等于投稿终稿，完整结果解读和中英文稿更新为下一步工作。
\end{document}
